% SIMBIG 2026 — Springer LLNCS format
\documentclass[runningheads]{llncs}

\usepackage[T1]{fontenc}
\usepackage{graphicx}
\usepackage{booktabs}
\usepackage{array}
\usepackage{multirow}
\usepackage{hyperref}
\usepackage{xcolor}
\usepackage{float}
% Reduce spacing around tables and figures
\setlength{\floatsep}{6pt plus 2pt minus 2pt}
\setlength{\textfloatsep}{8pt plus 2pt minus 4pt}
\setlength{\intextsep}{6pt plus 2pt minus 2pt}
\setlength{\abovecaptionskip}{4pt}
\setlength{\belowcaptionskip}{2pt}

\begin{document}

\title{Two-Stage Sexism Detection Pipeline for Spanish Twitter}

\titlerunning{Two-Stage Sexism Detection Pipeline}

\author{Gizela S. Thomas\inst{1} \and
 Sebasti\'{a}n Felipe Zambrano\inst{1}}

\authorrunning{G.S. Thomas and S.F. Zambrano}

\institute{IE University, Madrid, Spain \\
\email{gizela.thomas@gmail.com}}

\maketitle

%---------------------------------------------------------------
\begin{abstract}
Sexist content on social media disproportionately targets women and poses a growing
challenge for automated content moderation. This paper presents a two-stage
classification pipeline for detecting sexism in Spanish-language tweets, developed in
collaboration with the United Nations International Computing Centre (UNICC). Stage~1
performs binary classification on a unified corpus of 44,352 tweets drawn from four
publicly available datasets, and Stage~2
categorizes confirmed sexist tweets into one of four semantic classes. Four transformer
models were benchmarked: BETO, RoBERTuito, XLM-T, and BERTIN. BETO achieved
the strongest overall performance with a Stage~1 Macro F1 of 0.7856 and Stage~2 Macro
F1 of 0.5852. The tight clustering of results across all models points to a data quality
ceiling rather than an architectural limitation, and the 0.20 point drop from Stage~1 to
Stage~2 reflects structural class imbalance rather than model failure. Fine-tuned
open-source models prove competitive with prompt-engineered large language model
approaches at zero inference cost, supporting their viability for resource-constrained
moderation systems such as the UNICC.

\keywords{Sexism detection \and Spanish NLP \and Text classification \and
Transformer models \and Content moderation \and Hate speech}
\end{abstract}

%---------------------------------------------------------------
\section{Introduction}

One in three women experiences some form of sexual harassment online, making social
media platforms an increasingly hostile space for gender-based discourse. As daily
communication increasingly takes place online, social media platforms have also become
an environment where sexist and hateful content directed at women continues to spread.
Automated detection systems powered by natural language processing offer a scalable
solution for identifying and moderating such content before it causes harm.

This paper presents a two-stage classification pipeline for detecting sexism in
Spanish-language tweets, developed in collaboration with the United Nations
International Computing Centre (UNICC). The models were trained on a unified corpus
of 44,352 tweets, combining four publicly available datasets and a refined taxonomy
that consolidates 16 original subcategories into 4 linguistically defensible classes.
Stage~1 performs binary classification to determine whether a tweet contains sexist
content. Stage~2 categorizes confirmed sexist tweets into one of four semantic
categories. We benchmark four Spanish and Twitter-specific pre-trained language
models --- BETO, RoBERTuito, XLM-T, and BERTIN --- and compare results against
existing approaches, including prompt engineering methods using large language models,
to assess the viability of fine-tuned models for low-cost, deployable content moderation
systems.

%---------------------------------------------------------------
\section{Related Work}

The United Nations has identified the elimination of violence against women and the
promotion of technology for gender equality as core development priorities. In pursuit
of these goals, advanced natural language processing techniques have been applied to
the automatic detection of sexist content online. In collaboration with UNICC, several
research teams have developed machine-learning and large-language-model approaches
to classify sexist tweets using a shared dataset. The Universitat Polit\`ecnica de
Val\`encia applied transformer-based models, including BETO and RoBERTa, to detect
sexist content in Spanish-language tweets. IE University extended this work by
introducing multi-class analysis, identifying five semantic categories to capture nuances
in the type of sexism expressed. Most recently, the Universidade da Coru\~na explored
prompt engineering techniques using large language models for the same classification
task~\cite{limaylla2025}. The present study builds on this line of work by introducing
a two-stage pipeline that combines binary detection with multi-class categorization and
benchmarks four Spanish and Twitter-specific fine-tuned models against this task.

%---------------------------------------------------------------
\section{Dataset}

\subsection{Source Datasets}

Table~\ref{tab:datasets} summarizes the core datasets used as the foundation for model
training and evaluation.

\begin{table}[h]
\caption{Summary of core datasets used for sexism and hate speech detection.}
\label{tab:datasets}
\centering
\begin{tabular}{lllll}
\toprule
\textbf{Feature} & \textbf{EXIST 2021} & \textbf{HatEval 2019} & \textbf{MeToo} & \textbf{EDOS} \\
\midrule
Language(s) & EN, ES & EN, ES & ES & EN \\
Total size & 6,930 & 15,126 & 2,282 & 20,016 \\
Label type & Binary & Binary + fine-grained & Binary & Binary + multi-class \\
\bottomrule
\end{tabular}
\end{table}

The EXIST~2021 dataset contains approximately 7,000 English and Spanish tweets
categorized as sexist or non-sexist. The ManRosSexism\_Twitter\_MeToo dataset,
compiled during the Me~Too movement, consists of 2,282 Spanish-language tweets
labeled as sexist or non-sexist. The HatEval dataset, released as part of SemEval~2019,
contains approximately 15,000 combined English and Spanish tweets annotated for hate
speech targeting women and immigrants. The EDOS dataset contains approximately
20,000 English tweets with binary labels, with a portion containing fine-grained
multi-class annotations. The combined unified dataset totals 44,352 tweets.

\subsection{Dataset Unification}

The dataset used for this study is taken from IE University's research team, which
translated all tweets into Spanish using Google Translate in Google Sheets. The raw
data was distributed across multiple CSV files. Using a Python preprocessing pipeline
in Google Colab, the datasets were consolidated by loading and concatenating files,
removing duplicates, standardizing label formats, dropping incomplete entries, and
filtering off-topic content. Each dataset was normalized to a consistent binary label
(1~=~sexist, 0~=~non-sexist), with subcategory annotations preserved where available.

\subsection{Taxonomy Consolidation}

The original EDOS dataset included 11 subcategories of sexist content with overlapping
definitions, making consistent annotation and modeling difficult. We manually reviewed
each subcategory and consolidated them into 5 broader classes based on tone, intent,
and target: Sexual Insults \& Objectification, Gendered Stereotypes \& Insults, General
Hostile Language, Threats \& Physical Harm, and Victim Blaming \& Justification.

\begin{figure}[H]
\centering
\includegraphics[width=0.65\textwidth]{class_distribution.png}
\caption{Distribution of multi-class sexism categories (n~=~8,612).}
\label{fig:dist}
\end{figure}

Upon reviewing the class distribution, Victim Blaming \& Justification contained only
107 examples, a significant imbalance relative to the majority class. Given the
insufficient volume of training examples and its linguistic proximity to Gendered
Stereotypes \& Insults --- both categories involve the justification or minimization of
harm directed at women --- we merged the two into a single class. This produced a
final 4-class taxonomy totaling 8,612 labeled examples (Fig.~\ref{fig:dist}). General
Hostile Language accounts for approximately 77\% of the multi-class subset, a 15:1
imbalance relative to minority classes, addressed through weighted cross-entropy loss
during training (Section~\ref{sec:training}).

During the review process, two subcategories labeled ``untargeted~| non-aggressive''
and ``untargeted~| aggressive'' were flagged as containing predominantly anti-immigrant
or racially motivated speech, derived from the HatEval dataset, which also examined
hate speech directed at immigrant groups.

%---------------------------------------------------------------
\section{Methodology}

\subsection{Two-Stage Pipeline Architecture}

This study introduces a two-stage classification pipeline for sexism detection in
Spanish-language tweets. Table~\ref{tab:pipeline} summarizes the key differences
between the two stages.

\begin{table}[h]
\caption{Key differences between Stage~1 and Stage~2 of the pipeline.}
\label{tab:pipeline}
\centering
\begin{tabular}{lll}
\toprule
\textbf{Feature} & \textbf{Stage 1} & \textbf{Stage 2} \\
\midrule
Task & Binary classification & Multi-class classification \\
Training data & 44,352 tweets (full dataset) & Sexist tweets only \\
Output & Sexist / Not sexist & 4 semantic categories \\
Purpose & Filter non-sexist content & Categorize type of sexism \\
\bottomrule
\end{tabular}
\end{table}

During inference, a tweet is first passed through Stage~1; if classified as non-sexist,
the pipeline terminates. If classified as sexist, it is forwarded to Stage~2 for category
assignment, as illustrated in Fig.~\ref{fig:pipeline}.

\begin{figure}[H]
\centering
\includegraphics[width=0.78\textwidth]{pipeline.png}
\caption{Two-stage sexism detection pipeline.}
\label{fig:pipeline}
\end{figure}

This architecture mirrors
real-world content moderation workflows, where a broad detection layer precedes more
granular analysis. Key advantages include modularity (each stage can be fine-tuned
independently), efficiency (non-sexist tweets never reach Stage~2), and task alignment
(each model trains only on relevant data).

\subsection{Models}

Four pre-trained transformer models were evaluated across both stages, selected to
represent a range of pretraining approaches from general Spanish corpora to
Twitter-specific data. Table~\ref{tab:models} summarizes their key characteristics.

\begin{table}[h]
\caption{Pre-trained models evaluated in this study.}
\label{tab:models}
\centering
\begin{tabular}{lllll}
\toprule
\textbf{Model} & \textbf{Arch.} & \textbf{Pretraining corpus} & \textbf{Lang.} & \textbf{Domain} \\
\midrule
BETO & BERT & Large Spanish corpus & ES & General \\
BERTIN & RoBERTa & Spanish web \& books & ES & General \\
XLM-T & XLM-R & 198M multilingual tweets & Multi & Twitter \\
RoBERTuito & RoBERTa & 500M Spanish tweets & ES & Twitter \\
\bottomrule
\end{tabular}
\end{table}

These four models were selected to enable a direct comparison of language specificity
versus domain alignment. Two additional models, XLM-RoBERTa-large and
mDeBERTa-v3, were included in the original benchmark design but could not be fully
evaluated due to hardware constraints, as discussed in Section~\ref{sec:limitations}.

\subsection{Training Setup}
\label{sec:training}

All code and notebooks are publicly available~\cite{thomas2026}. All models were fine-tuned using the HuggingFace Transformers library with PyTorch,
executed on an NVIDIA Tesla T4 GPU (14.56~GiB VRAM) via Google Colab.
Hyperparameters were selected based on established best practices for fine-tuning
transformer models for text classification, drawing on recommendations from the
HuggingFace documentation and prior work in Spanish NLP, and validated through
preliminary runs on the development set. Input sequences were truncated or padded to
128 tokens. Models were trained for up to 4 epochs with a batch size of 16, learning
rate of 2e-5, weight decay of 0.01, and warmup ratio of 0.1. Weighted cross-entropy
loss was applied with class weights computed inversely proportional to class frequency.
Early stopping with patience of 2 epochs optimized for Macro F1 on the validation set.
The dataset was split 70/15/15 (train/val/test), stratified by label.

\subsection{Evaluation Metrics}

The primary evaluation metric is Macro F1, which computes the unweighted average
F1 score across all classes, treating each class equally regardless of size. This is
appropriate for imbalanced datasets where accuracy would obscure poor performance
on minority categories. Accuracy, Weighted F1, Precision, and Recall are also reported.

%---------------------------------------------------------------
\section{Results}

Results are reported on the held-out test set (15\% of each stage's dataset:
approximately 6,653 tweets for Stage~1 and 1,292 for Stage~2), stratified by label.

\subsection{Stage 1 --- Binary Classification}

\begin{table}[h]
\caption{Stage~1 binary classification results. $\star$~=~best per metric.}
\label{tab:stage1}
\centering
\begin{tabular}{llllll}
\toprule
\textbf{Model} & \textbf{Accuracy} & \textbf{Macro F1} & \textbf{Weighted F1} & \textbf{Precision} & \textbf{Recall} \\
\midrule
BETO $\star$ & \textbf{0.8258} & \textbf{0.7856} & \textbf{0.8251} & \textbf{0.7880} & 0.7833 \\
RoBERTuito & 0.8247 & 0.7830 & 0.8235 & 0.7873 & 0.7791 \\
XLM-T & 0.8234 & 0.7829 & 0.8228 & 0.7849 & 0.7810 \\
BERTIN & 0.8130 & 0.7789 & 0.8159 & 0.7716 & \textbf{0.7887} \\
\bottomrule
\end{tabular}
\end{table}

All four models achieve similar Macro F1 scores ranging from 0.7789 (BERTIN) to
0.7856 (BETO), a spread of just 0.0067 points. BETO achieves the highest Macro F1
and accuracy, with its advantage likely stemming from robust general-purpose Spanish
representations spanning formal, colloquial, and regionally varied text. RoBERTuito and
XLM-T trail by a negligible margin (0.7830 and 0.7829 respectively), performing nearly
identically despite fundamentally different pretraining strategies.

BERTIN achieves the lowest Macro F1 (0.7789) but the highest recall (0.7887), meaning
it misses fewer sexist tweets at the cost of a slightly elevated false-positive rate
(precision: 0.7716). In a content moderation context where missing a sexist tweet
carries a higher social cost than over-flagging a borderline case for human review,
recall is arguably the more operationally important metric at this stage, making BERTIN
a strong candidate for deployment despite its lower overall Macro F1.

Macro F1 scores above 0.78 are competitive with prior work on the same dataset,
including prompt-engineering approaches using GPT-4o from the Universidade da
Coru\~na~\cite{limaylla2025}, transformer-based fine-tuning from the Universitat
Polit\`ecnica de Val\`encia, and multi-class baselines from IE University, suggesting
the field is converging on a similar performance range regardless of approach. The
similarity across four architecturally distinct models points to a data-quality bottleneck
rather than a model-selection problem. Further gains in Stage~1 are more likely to come
from dataset refinement than architectural improvements.

\subsection{Stage 2 --- Multi-class Classification}

Stage~2 categorizes sexist tweets into one of four semantic classes, trained on the full multi-class subset of 8,612 labeled examples. Table~\ref{tab:stage2} presents results across all four models.

\begin{table}[h]
\caption{Stage~2 multi-class classification results. $\star$~=~best per metric.}
\label{tab:stage2}
\centering
\begin{tabular}{llllll}
\toprule
\textbf{Model} & \textbf{Accuracy} & \textbf{Macro F1} & \textbf{Weighted F1} & \textbf{Precision} & \textbf{Recall} \\
\midrule
BETO $\star$ & \textbf{0.7933} & \textbf{0.5852} & \textbf{0.8045} & \textbf{0.5697} & 0.6128 \\
RoBERTuito & 0.7717 & 0.5762 & 0.7889 & 0.5464 & \textbf{0.6324} \\
XLM-T & 0.7693 & 0.5633 & 0.7842 & 0.5367 & 0.6115 \\
BERTIN & 0.7918 & 0.5526 & 0.7940 & 0.5738 & 0.5485 \\
\bottomrule
\end{tabular}
\end{table}

Macro F1 scores drop to a range of 0.5526--0.5852, roughly 0.20 points below Stage~1
across all models, though the ranking stays the same. The drop was expected given the
severity of class imbalance: General Hostile Language accounts for 77\% of the labeled
subset, a 15:1 ratio against minority classes, and Stage~2 trains on roughly one-tenth
of Stage~1's data while asking models to distinguish between categories that each
represent only 5--12\% of examples.

The semantic boundaries between categories also make Stage~2 genuinely difficult to
learn. A tweet like ``Porque las mujeres mienten y los hombres no. As\'i de simple''
sits between Gendered Stereotypes and General Hostile Language, while ``Lo que
probablemente caus\'o la violaci\'on en primer lugar, por lo que luego deben apedrearla
hasta la muerte'' falls between Sexual Insults and Threats. Higher recall from RoBERTuito
means more of these hard examples get correctly surfaced rather than absorbed into
General Hostile Language, which matters for any system that routes flagged content by
severity.

\subsection{Error Analysis}

This section examines systematic failure patterns across both stages using confusion
matrices and per-class accuracy analysis from BETO test set results.

\subsubsection{Stage 1 --- Binary Classification Errors.}

\begin{figure}[H]
\centering
\includegraphics[width=0.72\textwidth]{confusion_matrix_stage1.png}
\caption{Stage~1 confusion matrix --- BETO (dccuchile/bert-base-spanish-wwm-cased).}
\label{fig:cm1}
\end{figure}

Of the 6,653 test examples, BETO correctly classifies 3,949 non-sexist and 1,449 sexist
tweets, with a false positive rate of 16.72\% and a miss rate of 24.18\%. The 462 false
negatives are the most consequential error type: any sexist tweet that Stage~1 fails to
catch exits the pipeline entirely with no moderation flag. Several missed tweets
reference sexual harm in ways that read more like reporting than expressing sexism.
Others embed gendered insults in casual register, for example ``Callate que yo tmb te
extra\~no perra'', where hostility is softened by informal tone. These cases point to a
mislabeling problem as much as a modeling one. BERTIN's recall of 0.7887 versus
BETO's 0.7833 means fewer sexist tweets slip through undetected, which in a
harm-sensitive pipeline matters more than avoiding the occasional false alarm.

\subsubsection{Stage 2 --- Multi-class Classification Errors.}

\begin{figure}[H]
\centering
\includegraphics[width=0.72\textwidth]{confusion_matrix_stage2.png}
\caption{Stage~2 confusion matrix --- BETO (dccuchile/bert-base-spanish-wwm-cased).}
\label{fig:cm2}
\end{figure}

General Hostile Language achieves a class accuracy of 0.856 (854/998), while minority
classes trail substantially: Gendered Stereotypes \& Insults at 0.708 (114/161), Threats
\& Physical Harm at 0.591 (39/66), and Sexual Insults \& Objectification at just 0.313
(21/67). When uncertain, the model defaults to the majority class: 27 Gendered
Stereotypes, 19 Sexual Insults, and 15 Threats examples were all incorrectly predicted
as General Hostile Language.

The most striking failure is Sexual Insults \& Objectification, where the model gets more
predictions wrong than right (25 as Gendered Stereotypes, 19 as General Hostile
Language, versus 21 correct). Tweets that simultaneously sexualize and demean women
often straddle both categories, and short informal text without broader context makes
that distinction hard even for a human annotator. Improving Stage~2 will require more
than a better model: re-annotation by specialists in gender-based violence and hate
speech research, combined with targeted data collection for minority categories, would
likely have a greater impact than any architectural change.

%---------------------------------------------------------------
\section{Limitations}
\label{sec:limitations}

The class imbalance in the multi-class subset is an inherent characteristic of the source
dataset. General Hostile Language accounts for 77\% of the labeled subset, placing a
ceiling on Macro F1 performance regardless of model choice. Weighted cross-entropy
loss partially addresses this but does not fully compensate. English-language tweets were
machine-translated using Google Translate, which may introduce artifacts affecting model
learning. Manual inspection identified 133 tweets in the General Hostile Language
category as primarily non-gendered hate speech (4.6\% of that category), which may
slightly inflate performance metrics and highlights the challenge of intersectional hate
speech in dataset construction.

Two models could not be fully evaluated due to hardware constraints: XLM-RoBERTa-large
Stage~2 and mDeBERTa-v3 were excluded due to GPU memory limitations on the T4
used for training. The training data was predominantly sourced from Latin American
social media, while BETO and BERTIN were pretrained on corpora from Spain, which
may reduce effectiveness for regional expressions and slang. All results are based on a
single random seed without multiple runs; statistical significance testing is identified as
future work. Finally, the authors acknowledge a lack of formal training in gender or race
studies, which shaped how labels were interpreted during dataset construction.

%---------------------------------------------------------------
\section{Conclusions and Future Work}

This study presents a two-stage classification pipeline for sexism detection in
Spanish-language tweets, combining binary detection with fine-grained multi-class
categorization. BETO emerged as the strongest overall model, while RoBERTuito is
better suited for recall-sensitive deployment where missing a high-severity tweet carries
a higher cost than a false alarm. All four Stage~1 models clustered within 0.007 Macro
F1, pointing to a data quality ceiling. Spanish-specific models performed comparably to
or better than Twitter-domain XLM-T, suggesting language specificity is at least as
important as domain alignment. Fine-tuned open-source models proved competitive with
paid LLM approaches at zero inference cost. The 0.20 point Macro F1 drop from Stage~1
to Stage~2 reflects structural class imbalance rather than model failure; improving
Stage~2 requires better data, not better models. Manual inspection identified 4.6\% of
General Hostile Language as primarily non-gendered hate speech, surfacing
intersectionality as an unresolved challenge in dataset design.

Future work should prioritize a re-annotation effort led by specialists in gender studies,
hate speech research, and regional Spanish linguistics to improve label consistency and
reduce boundary ambiguity. Targeted data collection for minority categories, combined
with removal of machine-translation artifacts and non-gendered hate speech, would
address the core limitations more effectively than any modeling change. On the modeling
side, evaluating larger models on higher-memory hardware and applying data
augmentation for minority categories are natural next steps. Extending Stage~1 to a
three-way classifier separating sexist content, non-gendered hate speech, and neutral
content would directly address the intersectionality gap. Finally, incorporating computer
vision to detect sexism in images and memes represents an important extension, given
that harmful content online increasingly extends beyond text.

%---------------------------------------------------------------
\begin{credits}
\subsubsection{\ackname}
This work was developed in collaboration with the United Nations International
Computing Centre (UNICC) and builds on prior research conducted at IE University.

\subsubsection{\discintname}
The authors have no competing interests to declare.
\end{credits}

%---------------------------------------------------------------
\begin{thebibliography}{8}

\bibitem{limaylla2025}
Limaylla-Lunarejo, M.I.: Prompt Engineering in Spanish for Improving Sexism Detection
in Tweets. In: Proceedings of SIMBIG 2025. Springer (2025)

\bibitem{beto}
Ca\~nete, J., Chaperon, G., Fuentes, R., Ho, J.H., Kang, H., P\'erez, J.:
Spanish Pre-Trained BERT Model and Evaluation Data. In: PML4DC at ICLR 2020 (2020)

\bibitem{robertuito}
P\'erez, J.M., Furman, D.A., Alonso Alemany, L., Luque, F.: RoBERTuito:
A Pre-trained Language Model for Social Media Spanish. In: Proceedings of LREC 2022 (2022)

\bibitem{xlmt}
Barbieri, F., Camacho-Collados, J., Neves, L., Espinosa-Anke, L.:
TweetEval: Unified Benchmark and Comparative Evaluation for Tweet Classification.
In: Findings of EMNLP 2020 (2020)

\bibitem{bertin}
De la Rosa, J., Huertas-Tato, J., Gonz\'alez-Agirre, A., Villegas, M.:
BERTIN: Efficient Pre-Training of a Spanish Language Model. In: IberLEF 2021 (2021)

\bibitem{edos}
Kirk, H.R., Yin, W., Vidgen, B., Rowan, Z.: SemEval-2023 Task 10: Explainability
for Sexism in Online Discussions. In: Proceedings of SemEval 2023 (2023)

\bibitem{hateval}
Basile, V., Bosco, C., Fersini, E., et al.: SemEval-2019 Task 5: Multilingual Detection
of Hate Speech Against Immigrants and Women in Twitter. In: Proceedings of SemEval
2019 (2019)

\bibitem{exist2021}
Rodr\'iguez-S\'anchez, F., Carrillo-de-Albornoz, J., Plaza, L., et al.:
Overview of EXIST 2021 -- sEXism Identification in Social neTworks. In:
Proceedings of IberLEF 2021 (2021)

\bibitem{upv2023}
UN-ICC and Universitat Polit`ecnica de Val`encia: Women Safer Online -- Sexism Detection Pipeline. GitHub Repository (2023). \url{https://github.com/UN-ICC/UPV-capstone-Women-Safter-Online}

\bibitem{ie2024}
Pablo et al.: UNICC Capstone Project -- Sexism Detection. IE University GitHub Repository (2024). \url{https://github.com/pablo60217/capstoneprojectunicc}

\bibitem{thomas2026}
Thomas, G.S., Zambrano, S.F.: Two-Stage Sexism Detection Pipeline for Spanish Twitter. GitHub Repository (2026). \url{https://github.com/gizelat/sexist-language-twitter-spanish}

\end{thebibliography}

\end{document}
